\subsubsection{Diâmetro ($d_1 = \SI{90}{\milli\meter}$)}

\begin{itemize}
  \item \textbf{Flexão}
\end{itemize}

A tensão de flexão é dada por:

\begin{equation}
  \sigma_f = \frac{M_{f_e}}{W_f}
  \label{eq:sigma_f_d1}
\end{equation}

Onde:

\begin{itemize}
  \item $M_{f_e}$: momento fletor no eixo em \si{\newton\meter} (Tabela \ref{tab:dados_esforcos_tambor}).
  \item $W_f$: módulo de resistência à flexão em \si{\meter^3}.
\end{itemize}

\vspace{1em}

O módulo de resistência à flexão é dado por:

\begin{equation}
  W_f = \frac{\pi \cdot {d_1}^3}{32}
  \label{eq:w_f_d1}
\end{equation}

Onde:

\begin{itemize}
  \item $d_1$: diâmetro do eixo em \si{\milli\meter} \eqref{eq:d_1_eixo_do_tambor}.
\end{itemize}

\vspace{1em}

Substituindo os valores na equação \eqref{eq:w_f_d1}:

\begin{align*}
  W_f & = \frac{\pi \cdot {\left( 90 \cdot 10^{-3} \right)}^3}{32} \\
  W_f & = \boxed{\SI{7.16e-5}{\meter^3}}
\end{align*}

Substituindo os valores na equação \eqref{eq:sigma_f_d1}:

\begin{align*}
  \sigma_f & = \frac{\num{2785.41}}{\num{7.16e-5}} \\
  \sigma_f & = \boxed{\SI{38.9}{\mega\pascal}}
\end{align*}

\begin{itemize}
  \item \textbf{Torção}
\end{itemize}

A tensão de cisalhamento por torção é dada por:

\begin{equation}
  \tau_T = \frac{T_s}{W_T}
  \label{eq:tau_T_d1}
\end{equation}

Onde:

\begin{itemize}
  \item $T_s$: torque de saída do redutor em \si{\newton\meter} \eqref{eq:torque_saida_redutor}.
  \item $W_T$: módulo de resistência à torção em \si{\meter^3}.
\end{itemize}

\vspace{1em}

O módulo de resistência à torção é dado por:

\begin{equation}
  W_T = \frac{\pi \cdot {d_1}^3}{16}
  \label{eq:w_T_d1}
\end{equation}

Onde:

\begin{itemize}
  \item $d_1$: diâmetro do eixo em \si{\milli\meter} \eqref{eq:d_1_eixo_do_tambor}.
\end{itemize}

\vspace{1em}

Substituindo os valores na equação \eqref{eq:w_T_d1}:

\begin{align*}
  W_T & = \frac{\pi \cdot {\left( 90 \cdot 10^{-3} \right)}^3}{16} \\
  W_T & = \boxed{\SI{1.43e-4}{\meter^3}}
\end{align*}

Substituindo os valores na equação \eqref{eq:tau_T_d1}:

\begin{align*}
  \tau_T & = \frac{\num{9150.44}}{\num{1.43e-4}} \\
  \tau_T & = \boxed{\SI{63.99}{\mega\pascal}}
\end{align*}

\begin{itemize}
  \item \textbf{Cisalhamento}
\end{itemize}

A tensão de cisalhamento cortante é dada por:

\begin{equation}
  \tau_V = \frac{4}{3} \cdot \frac{V_e}{A}
  \label{eq:tau_V_d1}
\end{equation}

Onde:

\begin{itemize}
  \item $V_e$: força cortante no eixo em \si{\newton} (Tabela \ref{tab:dados_esforcos_tambor}).
  \item $A$: área da seção transversal do eixo em \si{\meter^2}.
\end{itemize}

\vspace{1em}

A área da seção transversal do eixo é dada por:

\begin{equation}
  A = \frac{\pi \cdot {d_1}^2}{4}
  \label{eq:a_d1}
\end{equation}

Onde:

\begin{itemize}
  \item $d_1$: diâmetro do eixo em \si{\milli\meter} \eqref{eq:d_1_eixo_do_tambor}.
\end{itemize}

\vspace{1em}

Substituindo os valores na equação \eqref{eq:a_d1}:

\begin{align*}
  A & = \frac{\pi \cdot {\left( 90 \cdot 10^{-3} \right)}^2}{4} \\
  A & = \boxed{\SI{6.36e-3}{\meter^2}}
\end{align*}

Substituindo os valores na equação \eqref{eq:tau_V_d1}:

\begin{align*}
  \tau_V & = \frac{4}{3} \cdot \frac{\num{27682.54}}{\num{6.36e-3}} \\
  \tau_V & = \boxed{\SI{5.8}{\mega\pascal}}
\end{align*}

\begin{itemize}
  \item \textbf{Tensão combinada}
\end{itemize}

A tensão combinada é dada por:

\begin{equation}
  \sigma_c = \sqrt{{\sigma_f}^2 + 3 \left( \tau_T + \tau_V \right)^2}
  \label{eq:sigma_c_d1}
\end{equation}

Substituindo os valores na equação \eqref{eq:sigma_c_d1}:

\begin{align*}
  \sigma_c & = \sqrt{\num{38.9}^2 + 3 \left( \num{63.99} + \num{5.8} \right)^2} \\
  \sigma_c & = \boxed{\SI{126.99}{\mega\pascal}}
\end{align*}

Verificando a tensão combinada com a tensão admissível:

\begin{align*}
  \sigma_c                  & \leq \sigma_{\text{adm}}                             \\
  \SI{126.99}{\mega\pascal} & \leq \SI{214.29}{\mega\pascal} \rightarrow \text{OK}
\end{align*}

\subsubsection{Diâmetro do assento ($d = \SI{60}{\milli\meter}$)}

\begin{itemize}
  \item \textbf{Torção}
\end{itemize}

A tensão de cisalhamento por torção é dada por:

\begin{equation}
  \tau_T = \frac{T_s}{W_T}
  \label{eq:tau_T_d}
\end{equation}

Onde:

\begin{itemize}
  \item $T_s$: torque de saída do redutor em \si{\newton\meter} \eqref{eq:torque_saida_redutor}.
  \item $W_T$: módulo de resistência à torção em \si{\meter^3}.
\end{itemize}

\vspace{1em}

O módulo de resistência à torção é dado por:

\begin{equation}
  W_T = \frac{\pi \cdot {d}^3}{16}
  \label{eq:w_T_d}
\end{equation}

Onde:

\begin{itemize}
  \item $d$: diâmetro do assento em \si{\milli\meter} \eqref{eq:d_eixo_do_tambor}.
\end{itemize}

\vspace{1em}

Substituindo os valores na equação \eqref{eq:w_T_d}:

\begin{align*}
  W_T & = \frac{\pi \cdot {\left( 60 \cdot 10^{-3} \right)}^3}{16} \\
  W_T & = \boxed{\SI{4.24e-5}{\meter^3}}
\end{align*}

Substituindo os valores na equação \eqref{eq:tau_T_d}:

\begin{align*}
  \tau_T & = \frac{\num{9150.44}}{\num{4.24e-5}} \\
  \tau_T & = \boxed{\SI{215.81}{\mega\pascal}}
\end{align*}

\begin{itemize}
  \item \textbf{Cisalhamento}
\end{itemize}

A tensão de cisalhamento cortante é dada por:

\begin{equation}
  \tau_V = \frac{4}{3} \cdot \frac{V_e}{A}
  \label{eq:tau_V_d}
\end{equation}

Onde:

\begin{itemize}
  \item $V_e$: força cortante no eixo em \si{\newton} (Tabela \ref{tab:dados_esforcos_tambor}).
  \item $A$: área da seção transversal do eixo em \si{\meter^2}.
\end{itemize}

\vspace{1em}

A área da seção transversal do eixo é dada por:

\begin{equation}
  A = \frac{\pi \cdot {d}^2}{4}
  \label{eq:a_d}
\end{equation}

Onde:

\begin{itemize}
  \item $d$: diâmetro do assento em \si{\milli\meter} \eqref{eq:d_eixo_do_tambor}.
\end{itemize}

\vspace{1em}

Substituindo os valores na equação \eqref{eq:a_d}:

\begin{align*}
  A & = \frac{\pi \cdot {\left( 60 \cdot 10^{-3} \right)}^2}{4} \\
  A & = \boxed{\SI{2.83e-3}{\meter^2}}
\end{align*}

Substituindo os valores na equação \eqref{eq:tau_V_d}:

\begin{align*}
  \tau_V & = \frac{4}{3} \cdot \frac{\num{27682.54}}{\num{2.83e-3}} \\
  \tau_V & = \boxed{\SI{13.04}{\mega\pascal}}
\end{align*}

\begin{itemize}
  \item \textbf{Tensão combinada}
\end{itemize}

Verificando as tensões de cisalhamento com a tensão admissível:

\begin{align*}
  \tau_T + \tau_V                                      & \leq \tau_{\text{adm}}                                   \\
  \SI{215.81}{\mega\pascal} + \SI{13.04}{\mega\pascal} & > \SI{214.29}{\mega\pascal} \rightarrow \text{Reprovado}
\end{align*}

Devido às tensões de cisalhamento serem maiores que a tensão admissível, é necessário aumentar o diâmetro do eixo para $d_1 = \SI{100}{\milli\meter}$.

\begin{table}[H]
  \centering
  \caption{Diâmetros do eixo do tambor corrigidos}
  \label{tab:diametros_eixo_do_tambor_corrigidos}
  \begin{tabularx}{\textwidth}{>{\centering\arraybackslash}X>{\centering\arraybackslash}X>{\centering\arraybackslash}X}
    \toprule
    Descrição                        & Símbolo & Valor                  \\
    \midrule
    Diâmetro do eixo                 & $d_1$   & \SI{100}{\milli\meter} \\
    Diâmetro do encosto do rolamento & $d_e$   & \SI{80}{\milli\meter}  \\
    Diâmetro do assento              & $d$     & \SI{70}{\milli\meter}  \\
    \bottomrule
  \end{tabularx}
\end{table}

Recalculando as tensões para o novo diâmetro do eixo:

\subsubsection{Diâmetro ($d_1 = \SI{100}{\milli\meter}$)}

\begin{align*}
  \sigma_f & = \num{2785.41} \cdot \frac{32}{\pi \cdot \num{0.1}^3} \\
  \sigma_f & = \boxed{\SI{28.37}{\mega\pascal}}
\end{align*}

\begin{align*}
  \tau_T & = \num{9150.44} \cdot \frac{16}{\pi \cdot \num{0.1}^3} \\
  \tau_T & = \boxed{\SI{46.6}{\mega\pascal}}
\end{align*}

\begin{align*}
  \tau_V & = \frac{4}{3} \cdot \frac{\num{27682.54}}{\pi \cdot \num{0.1}^2} \\
  \tau_V & = \boxed{\SI{4.7}{\mega\pascal}}
\end{align*}

\begin{align*}
  \sigma_c & = \sqrt{\num{28.37}^2 + 3 \left( \num{46.6} + \num{4.7} \right)^2}                      \\
  \sigma_c & = \boxed{\SI{93.27}{\mega\pascal}} \leq \SI{214.29}{\mega\pascal} \rightarrow \text{OK}
\end{align*}

\subsubsection{Diâmetro do assento ($d = \SI{70}{\milli\meter}$)}

\begin{align*}
  \tau_T & = \num{9150.44} \cdot \frac{16}{\pi \cdot \num{0.07}^3} \\
  \tau_T & = \boxed{\SI{135.87}{\mega\pascal}}
\end{align*}

\begin{align*}
  \tau_V & = \frac{4}{3} \cdot \frac{\num{27682.54}}{\pi \cdot \num{0.07}^2} \\
  \tau_V & = \boxed{\SI{9.59}{\mega\pascal}}
\end{align*}

\begin{equation*}
  \tau_T + \tau_V = \boxed{\SI{145.46}{\mega\pascal}} \leq \SI{214.29}{\mega\pascal} \rightarrow \text{OK}
\end{equation*}
